% =========================================================
\section{Validación funcional por subsistema}
\label{sec:validacion}
% =========================================================

La \cref{sec:sistema} describe lo que el sistema hace. Esta sección
establece con qué evidencia se sostiene esa descripción.

La dificultad de validar un sistema de automatización de laboratorio es
que su comportamiento depende de tres instrumentos cuyo acceso es
limitado: el láser y el osciloscopio solo están disponibles durante las
sesiones reservadas, y una parte de las condiciones que interesa
verificar --- la pérdida de un dispositivo a mitad de secuencia, el
rebote de una lectura alrededor del umbral de disparo, la escritura de
un registro incompleto --- no puede provocarse a voluntad frente al
equipo sin arriesgar la muestra o el instrumento.

El criterio adoptado fue separar la validación en dos planos. Todo lo
que depende de la lógica del programa --- la evaluación de las
condiciones de disparo, la construcción del registro de sesión, la
detección de anomalías en los datos guardados, la propagación de
advertencias a la interfaz --- se valida sin instrumento, mediante
simulación y sesiones sintéticas. Solo se reserva para el laboratorio
aquello que depende de la respuesta efectiva del equipo. Esta
separación no es una concesión a la escasez de tiempo de laboratorio:
es lo que permite ejercitar condiciones de falla que en el equipo real
serían irreproducibles.

\subsection{Simulación de la secuencia de medición}
\label{sec:val_simulador}

El módulo \texttt{app/medicion/trigger.py} concentra la lógica que
decide cuándo abrir y cuándo cerrar una adquisición. En modo por
temperatura esa decisión depende de una magnitud que varía con el
tiempo, con ruido, y que puede cruzar el umbral varias veces antes de
estabilizarse. Es la porción del sistema donde un error resulta más
difícil de advertir: una ventana que se cierra antes de tiempo produce
un archivo que parece válido.

Para ejercitarla se construyó \texttt{simular\_barrido.py}, que alimenta
al módulo de disparo con perfiles de temperatura sintéticos y verifica
que las ventanas de adquisición resultantes correspondan a lo esperado.
La suite consta de nueve casos que cubren el enfriamiento monótono, el
cruce del umbral con ruido, el rebote alrededor de la temperatura
objetivo y la terminación anticipada de la secuencia. Los nueve casos
resultan aprobados.

Ese resultado admite una salvedad que conviene declarar. El caso
\texttt{no\_cierra\_durante\_rebote} verifica que la ventana de
adquisición no se fragmente cuando la lectura oscila alrededor del
umbral, y lo hace comprobando que la duración de la ventana
correspondiente al objetivo de \SI{45}{\celsius} se mantenga entre una
décima parte y diez veces el promedio de las ventanas de sus objetivos
vecinos, \SI{40}{\celsius} y \SI{50}{\celsius}. En la corrida
verificada la ventana del objetivo central dura \SI{140.30}{\second}
de tiempo simulado, frente a \SI{130.32}{\second} y \SI{90.00}{\second}
de las vecinas, lo que arroja un cociente de 1.27 sobre un intervalo
admisible de entre 0.1 y 10.

Dos aspectos limitan el alcance de esta comprobación. El primero es la
amplitud del intervalo: dos órdenes de magnitud detectan una
fragmentación grosera de la ventana, no una desviación moderada. El
segundo, más relevante, es la naturaleza de la referencia. El criterio
no contrasta la ventana contra una duración esperada, sino contra el
promedio de sus vecinas, y esas vecinas difieren entre sí en cerca de
un cuarenta por ciento. Se trata, por tanto, de una prueba de
consistencia interna: detecta que una ventana se desmarque de las
demás ---que es precisamente lo que el rebote provocaría---, pero no
verifica que ninguna de ellas dure lo que debe durar. Se reporta como
aprobado con esa reserva explícita.

Conviene añadir una observación de orden práctico. Este caso simula el
enfriamiento de la muestra desde \SI{62}{\celsius} hasta
\SI{30}{\celsius} con un factor de compresión temporal de veinte, lo
que se traduce en cerca de catorce minutos de ejecución real frente a
los segundos que consumen los demás. Una prueba con ese costo se corre
con menos frecuencia que las restantes, y esa asimetría conviene
tenerla presente al interpretar la vigencia del resultado.

\subsection{Validación de la cadena de medición y almacenamiento}
\label{sec:val_cadena}

La simulación anterior ejercita la lógica de disparo de forma aislada.
Un segundo nivel de la misma herramienta sustituye los componentes
simulados por las clases reales de \texttt{app/medicion/medicion.py} y
\texttt{app/almacenamiento/almacenamiento.py}, dejando fuera únicamente
la comunicación con los instrumentos.

En esa configuración se recorre la cadena completa: se resuelve la
condición de disparo, se solicita la captura, se convierten las
muestras a unidades físicas, se escriben el archivo \texttt{.npy} y la
fila correspondiente del registro de sesión, y se vuelve a leer el
resultado. La corrida produce un registro de siete filas cuyo contenido
coincide con lo esperado tanto en el número de campos como en los
valores derivados de los metadatos de adquisición.

Lo que este nivel valida no es la calidad de la señal --- no hay señal
--- sino la integridad del camino que va de la decisión de capturar al
archivo en disco. Es precisamente el tramo donde se originaron dos de
las anomalías documentadas en la \cref{sec:fallas}.

\subsection{Compatibilidad con registros anteriores}
\label{sec:val_compatibilidad}

El formato del registro de sesión descrito en la
\cref{sec:formato_datos} no fue el primero. Las sesiones anteriores a la
integración del módulo de temperatura se guardaron con un conjunto
menor de columnas, y algunas de ellas conservan valor documental.

La lectura de sesiones se implementó tratando el conjunto de campos
declarado en \texttt{COLUMNAS\_REQUERIDAS} como un subconjunto del
encabezado presente en el archivo, en lugar de exigir coincidencia
exacta. Con ese criterio, un registro de dieciocho columnas se abre sin
modificación y sin pérdida de los campos que sí contiene. La
verificación se realizó sobre sesiones reales conservadas del formato
anterior.

La contraparte de esa decisión es que la ausencia de una columna deja de
ser un error de apertura y pasa a ser una condición que debe detectarse
más tarde, en la verificación del contenido. Ese desplazamiento es
deliberado y se examina en la subsección siguiente.

\subsection{Verificación de sesiones y validación de la herramienta}
\label{sec:val_verificador}

La herramienta \texttt{verificar\_sesion.py} recorre una sesión ya
guardada y evalúa diez condiciones anómalas: ausencia de archivos de
muestras (\texttt{npy\_ausente}), imposibilidad de reconstruir captura
alguna (\texttt{sin\_reconstruibles}), saturación de la señal
(\texttt{saturacion}), repetición de un mismo contenido binario bajo
la misma escala o bajo escalas distintas
(\texttt{buffer\_repetido\_misma\_escala} y
\texttt{buffer\_repetido\_escala\_distinta}), coexistencia de escalas
diferentes en una misma serie (\texttt{mezcla\_escalas}), deriva del
tiempo de arribo (\texttt{deriva\_onset}), estancamiento de la lectura
de temperatura (\texttt{temperatura\_estancada}), presencia de
capturas marcadas con error (\texttt{error\_flag}) y dispersión
excesiva de la amplitud pico a pico (\texttt{dispersion\_vpp}). Es la
herramienta con la que se detectaron las anomalías descritas en la
\cref{sec:falla_buffer} y en la \cref{sec:falla_regresion}.

Una herramienta de auditoría que nunca se ha visto fallar no es una
herramienta verificada, sino una herramienta afortunada. Su validación
se realizó en dos etapas.

La primera consistió en fabricar, mediante
\texttt{generar\_sesion\_prueba.py}, once sesiones sintéticas: diez
diseñadas para provocar cada una de las condiciones anteriores, y una
undécima deliberadamente limpia. Esta última no es un caso menor. Un
verificador que declarara anomalías en toda sesión que examina
detectaría las diez condiciones sin dificultad y sería inútil; la
sesión limpia es la que exige que no las declare cuando no las hay. El
verificador identifica las diez condiciones y no emite alerta alguna
sobre la sesión limpia.

La segunda etapa responde a una objeción que la primera deja abierta.
El corredor compara las alertas obtenidas contra un conjunto esperado,
y ese conjunto podría haberse derivado de la salida observada del
verificador en lugar del comportamiento pretendido; de ser así, la
prueba pasaría igual con el verificador averiado. Para descartarlo se
realizó una prueba de mutación: se desactivó la emisión de cada una de
las diez condiciones, una a la vez, restaurando el archivo entre
mutaciones, y se registró qué casos dejaban de detectarse.

Las diez mutaciones provocaron la caída del caso correspondiente.
Ninguna sobrevivió, y en ningún caso el conjunto completo quedó en
verde con un chequeo desactivado. Dos de ellas arrastraron además un
caso ajeno: al desactivar \texttt{npy\_ausente} cae también
\texttt{sin\_reconstruibles}, y al desactivar \texttt{mezcla\_escalas}
cae \texttt{buffer\_repetido\_escala\_distinta}. Ambas muertes
cruzadas corresponden a relaciones de implicación entre las
condiciones --- una sesión sin capturas reconstruibles es
necesariamente una sesión a la que le faltan archivos de muestras, y
la repetición de contenido bajo escalas distintas presupone escalas
distintas ---, previstas por análisis del código antes de la prueba y
confirmadas empíricamente por ella.

La mutación dejó un hallazgo que no se buscaba. Al desactivar
\texttt{sin\_reconstruibles}, el código de salida del verificador
permanece en 1: la instrucción de retorno anticipada devuelve una lista
que \texttt{npy\_ausente} ya había poblado en la condición anterior. Un
corredor que juzgara el resultado por el código de salida habría dado
ese caso por aprobado con el chequeo apagado. Lo que impidió el falso
aprobado fue la decisión de comparar el conjunto de alertas emitidas
contra el esperado, y no únicamente el código de salida. La
observación no es anecdótica: reaparece en
\cref{sec:val_advertencias}, donde una comprobación distinta calcula su
código de retorno antes del paso que pretende verificar. En ambos
casos, el código de salida resultó ser el indicador menos informativo
disponible.

Queda una limitación que conviene declarar. La prueba de mutación se
aplicó manualmente y no existe en el repositorio como procedimiento
automatizado: si \texttt{verificar\_sesion.py} incorpora una condición
nueva o reescribe una existente, el resultado anterior deja de ser
válido y nada lo advierte. La herramienta que somete al sistema a
verificación continua está, ella misma, verificada por una prueba de
un solo momento.

\subsection{Propagación de advertencias a la interfaz}
\label{sec:val_advertencias}

El sistema no interrumpe una secuencia ante una anomalía recuperable:
la registra y la señala. Esa decisión, justificada en la
\cref{sec:decisiones}, hace que el indicador de advertencia de la
interfaz sea el único aviso que recibe el operador mientras la sesión
transcurre.

La comprobación \texttt{probar\_advertencias\_gui.py} ejercita cuatro
trayectorias de esa señalización y las cuatro resultan conformes. El
alcance de esa cifra, sin embargo, es menor de lo que sugiere: la
comprobación calcula su código de retorno antes de ejecutar el único
paso que requiere confirmación visual del operador, de modo que ese
paso no puede hacerla reprobar. Los cuatro resultados corresponden a
las trayectorias automatizables; el paso restante queda como
comprobación asistida, y así se reporta.

\subsection{Verificaciones con instrumento}
\label{sec:val_instrumento}

La separación adoptada deja fuera un residuo que no admite sustituto
por simulación. Cinco comprobaciones dependen de la respuesta efectiva
del equipo y quedaron declaradas como pendientes hasta la sesión de
laboratorio del 8 de septiembre de 2026, en la que se ejecutaron las
cinco. Se reportan aquí junto con el criterio de aceptación fijado de
antemano, porque dos de ellas no lo cumplieron en su primera ejecución
y esa discrepancia forma parte del resultado.

\subsubsection*{Poblado de las columnas de escala}

La regresión descrita en la \cref{sec:falla_regresion} se corrigió
restituyendo las consultas \texttt{CH1:SCALE?} y
\texttt{HORizontal:SCAle?} dentro de la rutina de lectura del
preámbulo. El criterio fijado excluía explícitamente el valor
\num{0.0} como aprobación, pues es lo que escriben los manejadores de
excepción cuando la consulta se lanza y falla: solo un número
consistente con el panel del instrumento confirma la corrección. En la
sesión \texttt{20260908\_184623} las once capturas registran
\texttt{CH\_SCALE} $=\num{0.001}$ y \texttt{HOR\_SCALE} $=\num{1e-5}$,
equivalentes a \SI{1}{\milli\volt} y \SI{10}{\micro\second} por
división, en coincidencia con el panel. Conforme.

\subsubsection*{Captura manual con el módulo de temperatura
desconectado}

El criterio fijado exigía tres cosas de la fila resultante:
\texttt{error\_flag = 1}, un \texttt{error\_desc} que nombrara al
dispositivo ausente y un campo de temperatura en cero. La primera
ejecución cumplió únicamente la primera. Las capturas
\texttt{m0012}--\texttt{m0014} de la sesión \texttt{20260908\_184623}
levantaron la bandera, pero escribieron \SI{24.25}{\celsius} sin sensor
conectado y una descripción genérica. El defecto y su corrección se
documentan en la \cref{sec:falla_fantasma}.

La corrección modificó además el propio criterio. El campo en cero se
sustituyó por \texttt{nan}, porque \SI{0}{\celsius} es una temperatura
físicamente válida y un análisis posterior no podría distinguirla de
una medición; \texttt{nan} no admite esa lectura. En la sesión
\texttt{20260908\_190614}, con el módulo desconectado, las filas
registran \texttt{error\_flag = 1}, \texttt{temperatura = nan} y un
\texttt{error\_desc} que nombra al ESP32 y su puerto. Conforme bajo el
criterio corregido.

\subsubsection*{Limpieza automática de la bandera del monitor}

El criterio exigía que la bandera del monitor de conexión se limpiara
por sí sola en el siguiente ciclo de consulta tras restablecerse el
módulo. La primera ejecución tampoco lo cumplió: el registro de la
sesión \texttt{20260908\_190614} muestra que cada recuperación
(\texttt{m0005} y \texttt{m0013}) ocurre después de una pausa
prolongada, la que toma accionar el botón de reconexión en la
interfaz. No había limpieza automática porque no había reconexión
automática: el método correspondiente se limitaba a informar el estado
de la conexión sin intentar restablecerla.

Se implementó el reintento sobre el puerto configurado y, al recibirse
una trama válida, el reinicio del hilo de lectura. La sesión
\texttt{20260908\_200938} lo verifica: las capturas \texttt{m0003} y
\texttt{m0004} se registran con el módulo desconectado y
\texttt{temperatura = nan}, y \texttt{m0005} vuelve a
\texttt{error\_flag = 0} con lectura real, sin intervención del
operador.

La implementación reintenta únicamente sobre el puerto declarado en la
configuración y no recorre los puertos disponibles buscando el
dispositivo. La restricción es deliberada. Un puerto no se reasigna
solo: cuando cambia es porque alguien manipuló la conexión física o el
sistema operativo la renumeró tras esa manipulación, y en ambos casos
hay una persona presente que puede corregir la configuración. Recorrer
los puertos implicaría, en cambio, abrir puertos ajenos al módulo
---entre ellos el del láser--- para averiguar qué hay detrás de cada
uno, lo que introduce un riesgo mayor que el problema que resuelve.

El costo de esa decisión se materializó en la misma sesión: el módulo
apareció en \texttt{COM5} al inicio y en \texttt{COM4} al final, y esa
transición sí requirió editar la configuración y reiniciar la
aplicación. La reconexión automática cubre la pérdida de comunicación
sobre un puerto estable, no la reasignación del puerto.

\subsubsection*{Recorrido completo del modo por temperatura}

La cuarta comprobación es la que ninguna simulación puede sustituir:
el modo de adquisición gobernado por temperatura, ejercitado de forma
exhaustiva según la \cref{sec:val_simulador}, no había recorrido nunca
de punta a punta los tres instrumentos simultáneos.

La sesión \texttt{20260908\_203856} lo recorrió. Con la muestra
calentada por encima del rango de interés y midiendo sobre su
enfriamiento pasivo, la secuencia produjo nueve puntos entre
\SI{43.89}{\celsius} y \SI{39.91}{\celsius}, sobre un barrido
programado de \SI{44}{\celsius} a \SI{40}{\celsius} en pasos de
\SI{0.5}{\celsius}, con \texttt{error\_flag = 0} en los nueve, las
columnas de escala pobladas y el campo \texttt{pulsos\_estimados}
---que contabiliza los disparos del láser integrados dentro de cada
ventana de adquisición--- con valor por primera vez fuera de
simulación. La lógica de apertura y cierre de ventana, verificada hasta
entonces contra perfiles sintéticos, gobernó una adquisición real.

La captura \texttt{m0001} del mismo archivo no pertenece al barrido:
es la adquisición manual de verificación que el sistema solicita antes
de iniciar cualquier secuencia, para que el operador confirme que la
señal es razonable y elija el destino de los datos. Su campo
\texttt{pulsos\_estimados} está vacío porque no hubo ventana de
acumulación que medir, y su modo registrado es \texttt{manual}. El
formato distingue ambas cosas sin ambigüedad.

Los objetivos efectivamente alcanzados fueron \SI{43.89}{\celsius},
\SI{43.39}{\celsius}, \SI{42.88}{\celsius}, \SI{42.39}{\celsius},
\SI{41.91}{\celsius}, \SI{41.41}{\celsius}, \SI{40.89}{\celsius},
\SI{40.39}{\celsius} y \SI{39.91}{\celsius}. La separación entre
puntos consecutivos se aparta del paso nominal de
\SI{0.5}{\celsius} en a lo sumo \SI{0.03}{\celsius}. El intervalo
entre capturas crece de forma monótona, de \SI{81}{\second} a
\SI{100}{\second}, lo que corresponde al enfriamiento pasivo
frenándose conforme la muestra se acerca a la temperatura ambiente:
el sistema no gobierna ese ritmo, solo espera a que el umbral se
cumpla.

\subsubsection*{Exportación al formato \texttt{.mat}}

La quinta comprobación se ejecutó sobre los datos de esa misma sesión.
El archivo resultante contiene el arreglo crudo, las muestras
convertidas a volts, el eje temporal reconstruido, la frecuencia de
muestreo, el número de puntos y los veintiún campos de metadatos del
registro de sesión. Conforme.

\subsubsection*{Lo que estas cinco comprobaciones dejan ver}

Tres de las cinco resultaron conformes en su primera ejecución. Dos no,
y ninguna de las dos podía haberse detectado sin el instrumento: una
escribía un dato inexistente y la otra ofrecía una función que nunca
llegó a implementarse. Ambas corresponden a trayectorias que la
validación por simulación no alcanza por construcción y no por descuido
en su diseño: \texttt{generar\_sesion\_prueba.py} fabrica el archivo
que produce una desconexión, no la desconexión misma.

El resultado respalda la separación adoptada al inicio de esta sección
y marca al mismo tiempo su frontera. La simulación verificó
correctamente todo lo que estaba a su alcance; lo que quedó fuera de su
alcance falló en dos de cinco casos.

\begin{table}[H]
  \centering
  \caption{Resumen de la validación por subsistema. Los casos marcados
           como asistidos requieren confirmación visual del operador.}
  \label{tab:resumen_validacion}
  \begin{tabular}{llcl}
    \toprule
    \textbf{Subsistema} & \textbf{Herramienta} & \textbf{Casos} &
    \textbf{Resultado} \\
    \midrule
    Lógica de disparo & \texttt{simular\_barrido.py} & 9 &
      Conforme\textsuperscript{a} \\
    Medición y almacenamiento & \texttt{simular\_barrido.py} (nivel 2) & 7 &
      Conforme \\
    Verificación de sesiones & \texttt{generar\_sesion\_prueba.py} &
      $10+1$\textsuperscript{c} & Conforme \\
    Advertencias en interfaz & \texttt{probar\_advertencias\_gui.py} & 4 &
      Conforme\textsuperscript{b} \\
    Verificaciones con instrumento & Sesión de laboratorio & 5 &
      Conforme\textsuperscript{d} \\
    \bottomrule
  \end{tabular}

  \vspace{0.5em}
  \begin{minipage}{0.85\textwidth}
    \footnotesize
    \textsuperscript{a} El caso \texttt{no\_cierra\_durante\_rebote}
    aprueba bajo un criterio de tolerancia amplio; véase la
    \cref{sec:val_simulador}.\\
    \textsuperscript{b} Cuatro trayectorias automatizables; un paso
    adicional queda como comprobación asistida.\\
    \textsuperscript{c} Diez condiciones de alerta más una sesión limpia
    de control negativo.\\
    \textsuperscript{d} Dos de las cinco exigieron corregir el sistema
    antes de resultar conformes, y una de ellas exigió además corregir el
    criterio; véase la \cref{sec:val_instrumento}.
  \end{minipage}
\end{table}
